\documentclass[12pt, french, latin.medieval, a4paper]{report}
% package de langues 
\usepackage[french]{babel}
\usepackage[utf8]{inputenc}
\usepackage[T1]{fontenc}

% package pour référence de liens sur pdf 
\usepackage[pdftex]{hyperref} % liens
\usepackage{lastpage} % numéro de page
\usepackage{caption}

% package de formes de document 
\usepackage[left=2.5cm,
              right=2.5cm,
              top=2.5cm,
              bottom=2.5cm
             ]{geometry}
% pour les tableaux 
\usepackage{array,multirow,makecell} 
\newcolumntype{R}[1]{>{\raggedleft\arraybackslash }b{#1}}
\newcolumntype{L}[1]{>{\raggedright\arraybackslash }b{#1}}
\newcolumntype{C}[1]{>{\centering\arraybackslash }b{#1}}

% package pour la page de couverture 
\usepackage{graphicx}
\usepackage[cyr]{aeguill}
\usepackage{fancyheadings}
\usepackage{amsmath}
\usepackage{amsfonts}
\usepackage{amssymb}
\usepackage{geometry}
\usepackage{eso-pic}
\usepackage{tikz}
\usetikzlibrary{calc}


% nouvelle commande pour la forme 
%\fancyhf[OFC]{\thepage / \pageref{LastPage}}
\rfoot{\thepage / \pageref{LastPage}}
%\setcounter{tocdepth}{3}

% pakages de couleurs du document
\usepackage{color}
\usepackage{xcolor}


% package pour la chimie
\usepackage[version=4]{mhchem}


% package pour la visualisation du code
\usepackage{listings}
\definecolor{darkWhite}{rgb}{0.94,0.94,0.94}
 
\lstset{
  aboveskip=3mm,
  belowskip=-2mm,
  backgroundcolor=\color{darkWhite},
  basicstyle=\footnotesize,
  breakatwhitespace=false,
  breaklines=true,
  captionpos=b,
  commentstyle=\color{red},
  deletekeywords={...},
  escapeinside={\%*}{*)},
  extendedchars=true,
  framexleftmargin=16pt,
  framextopmargin=3pt,
  framexbottommargin=6pt,
  frame=tb,
  keepspaces=true,
  keywordstyle=\color{blue},
  language=python,
  literate=
  {²}{{\textsuperscript{2}}}1
  {⁴}{{\textsuperscript{4}}}1
  {⁶}{{\textsuperscript{6}}}1
  {⁸}{{\textsuperscript{8}}}1
  {€}{{\euro{}}}1
  {é}{{\'e}}1
  {è}{{\`{e}}}1
  {ê}{{\^{e}}}1
  {ë}{{\¨{e}}}1
  {É}{{\'{E}}}1
  {Ê}{{\^{E}}}1
  {û}{{\^{u}}}1
  {ù}{{\`{u}}}1
  {â}{{\^{a}}}1
  {à}{{\`{a}}}1
  {á}{{\'{a}}}1
  {ã}{{\~{a}}}1
  {Á}{{\'{A}}}1
  {Â}{{\^{A}}}1
  {Ã}{{\~{A}}}1
  {ç}{{\c{c}}}1
  {Ç}{{\c{C}}}1
  {õ}{{\~{o}}}1
  {ó}{{\'{o}}}1
  {ô}{{\^{o}}}1
  {Õ}{{\~{O}}}1
  {Ó}{{\'{O}}}1
  {Ô}{{\^{O}}}1
  {î}{{\^{i}}}1
  {Î}{{\^{I}}}1
  {í}{{\'{i}}}1
  {Í}{{\~{Í}}}1,
  morekeywords={*,...},
  numbers=left,
  numbersep=10pt,
  numberstyle=\tiny\color{black},
  rulecolor=\color{black},
  showspaces=false,
  showstringspaces=false,
  showtabs=false,
  stepnumber=1,
  stringstyle=\color{gray},
  tabsize=4,
  title=\lstname,
}

%package pour les algorithmes
\usepackage{algorithm}
\usepackage{algorithmic}
\usepackage{amsmath}

% package pour les tableaux
\usepackage{longtable}
\usepackage{booktabs}

% package pour les figures 
\usepackage{subcaption}

% package de reférencement
% \usepackage[backend=biber,safeinputenc, sorting=none]{biblatex}

% mettre la bibliographie dans la table de matières
%\usepackage[nottoc, notlof, notlot]{tocbibind}
\include{page_couverture.tex}


\author{MBE MBE MINDJANA LOIC HENRI}
\title{GNPS et MDByTE}

\begin{document}
    
    \pagecouverture
    {A Sample-Centric and Knowledge-Driven Computational Framework for Natural
Products Drug Discovery}
    {}
    \tableofcontents
    \vspace{20cm}
    
    \section{Contexte et description de l'article}
    Dans le but d'aider à la recherche des médicaments basées sur les produits naturelles, plusieurs solutions récemment proposées utilisent des étapes d'alignements de fonctionnalités à l'instar du processus FBMN (Featured-based Molecular Networking) (présent sur la figure \ref{fig:fbmn}) qui rendent coûteux l'ajout de nouveaux exemples d'extraits.
    De plus, aucune d'elle ne propose de moyen d'analyse permettant de correspondre efficacement des exemples avec les données métabolomiques et métadonnées de sources externes.
    Dans l'objectif de répondre à ces limites les auteurs ont proposés dans l'article de référence \cite{gaudry2023sample}, un framework ENPKG(Experimental Natural Products Knowledge Graph)inspiré de LOTUS \cite{rutz2021lotus}, basé sur les graphes de connaissances et centré sur l'exemple, qui sera sujet de ce document.
 
    Comme informations liées à l'article, on a:
    \begin{itemize}
        \item \textbf{titre} :{A Sample-Centric and Knowledge-Driven Computational Framework for Natural Products Drug Discovery}
        \item \textbf{auteurs}: {Gaudry, Arnaud and Pagni, Marco and Mehl, Florence and Moretti, S{\'e}bastien and Quiros-Guerrero, Luis-Manuel and Rutz, Adriano and Kaiser, Marcel and Marcourt, Laurence and Queiroz, Emerson Ferreira and Ioset, Jean-Robert and others}
        \item \textbf{année}: {2023}
    \end{itemize}

    \begin{figure}
        \centering
        \includegraphics[width = \linewidth]{Images/fbmn.png}
        \caption{Workflow FBMN et MN classique \cite{nothias2020feature}}
        \label{fig:fbmn}
    \end{figure}

    \section{Workflow}
    Comme décrit plus haut, il était question pour les auteurs de proposer un framework qui permettrait de correspondre au mieux les exemples de leur système
    aux données et métadonnées métabolomiques de systèmes extérieurs. Le système obtenu dont workflow est présent sur la figure \ref{fig:enpkg}, permet d'introduire progressivement de nouveaux exemples.
    Pour etre ajouter dans le système, les exemples nécessite un formatage en dossier contenant :
    \begin{itemize}
        \item un dossier pour les fichiers pour le réseau moléculaire de l'extrait (de la liste de spectres MS2 de l'exemple)
        \item un dossier pour les notations "In Silico"
        \item un dossier pour la résolution de typologie de l'exemple sur le graphe de connaissance
        \item un dossier pour les annotations SIRIUS
    \end{itemize}
    
    Ces formats d'exemples seront introduits par la suite dans le graphe de connaissance globale et sauvegardés dans une base de données \textbf{zenodo}.
    Ainsi pour contenir le graphe de connaissance global, on utilise la base de données de graphe  \textbf{\href{https://graphdb.ontotext.com}{GRAPHDB}}
	qui pourra être sujet de requête pour obtenir des informations globales internes comme externes. Cela à cause du fait que le graphe intègre des liaisons aux plateformes publiques externes telles que GNPS et ChEMBL respectant une certaine structure.
    
    \begin{figure}
        \centering
        \includegraphics[width = \linewidth]{Images/enpkg.png}
        \caption{Workflow du framework ENPKG \cite{gaudry2023sample}}
        \label{fig:enpkg}
    \end{figure}

	On peut segmenter le workflow en trois grandes parties qui se suivent
    
	\subsection{Acquisition des données}
    	Cette étape consiste à obtenir les données spectrales MS2 par chromatographie liquide LC-MS/MS des différents extraits de base. Ces spectres par exemple seront par la suite introduits dans le dépôt MASSIVE et utilisés pour le graphe de connaissance après prétraitement.

	\subsection{Prétraitement des exemples}
    	Cette étape consiste à prétraiter les exemples pour être introduit dans le dépôt MASSIVE et utilisés pour être introduit dans le graphe de connaissance après intégration.
    	Comme décrit brièvement plus haut, les exemples sont transformés en des dossiers correspondants d'intégration après ajout de métadonnées.

	\subsection{Intégration des exemples}
	C'est en cette phase qu'on introduit les exemples de base ou de nouvelles dans le graphe globale de connaissance.  Les dossiers d'exemples sont analysés et attachés aux bases de données WIKIDATA de références et ChEBML de molécules bioactives. Puis de ces dossiers sont générés les sous graphes à intégrer dans le graphe globale après introduction dans la base “zenodo".
    
	\section{Conclusion}
	En somme, il était question de présenter, tel que nous avons perçu le workflow ENPKG donné par les auteurs qui pourrait être très utile à la découverte de médicaments.
    Ce framework permet d'exploiter les connaissances externes (données et métadonnées externes) au système pour accélérer l'exploration des extraits du système (en effectuant des requêtes sur la base de données de graphe globale) et d’ajouter efficacement des exemples futurs.
	Ainsi d'après le workflow, il pourrait être intéressant d'intégrer l'utilisation des données RMN (1D et 2D) à l'instar de construire des réseaux moléculaires unifiés MS-RMN pour chacun des exemples d'extraits.
    Malgré les objectifs de sa création atteints, il a un inconvéniant majeur qui est l’utilisation d’un langage d'interrogation qui n’est pas facilement accessible par les utilisateurs finaux que sont les biologistes.
    Cela permet de se pencher sur la question de comment rendre moins complexe l’accès aux données par un utilisateur final? Comment construire automatiquement les requêtes sur la base de demandes textuelles des utilisateurs finaux?

\bibliographystyle{plain}
    \bibliography{reference_GNPS_MADByTE}
\end{document}




